\section{Gap and contributions}
\label{sec:gap}

The reviewed literature establishes that no surrogate family dominates, that winners are predictable from problem
features, that within-family configuration matters, and that approximation accuracy is a weak proxy for decision
quality. It does not, to our knowledge, contain a study that jointly varies (i) surface class and landscape features,
(ii) dimension and sample budget, (iii) noise level and type, (iv) fidelity structure and cost ratio, (v) the pairing of
the sampling scout with the deliverable model, (vi) the upscaling from learner to optimizer-ready surrogate, including
compression, and (vii) the optimizer---and scores complete pipelines by true-function regret per unit simulation cost,
while explaining the persistent disagreement between engineering practice (Gaussian processes) and tabular machine
learning (tree ensembles). The study part of this paper does this on the Surjanovic--Bingham functions, the
multi-fidelity analytical benchmarks, and two real engineering datasets: coupled finite-element sizing of an inflatable
ramjet duct, which carries a closed-form low-fidelity model at every point, and historical aircraft performance data.
